% GENERATED FILE. Edit canonical lessons, then run npm run generate:books.
% canonical-source-hash: fnv1a64:5bf44efc1e0bd9c7
% canonical-lessons: GU-C218-ugvu, GU-C218-suryoday, GU-C218-athamvu, GU-C218-suryast, GU-C218-ajvalu, GU-R218-review-213-218

\chapter{Sunrise and sunset}
\label{ch:gu-sunrise-and-sunset}
\hlchaptermodality{218}

\begin{chapteropening}
\textbf{By the end of this chapter:} \emph{I can say when the sun will rise and set.}

Complete the last lesson of chapter 218: I can say when the sun will rise and set.
\end{chapteropening}

\section[\texorpdfstring{\gu{ઊગવું} (ūgvũ)}{ઊગવું (ūgvũ)}]{\gu{ઊગવું} (ūgvũ) — to rise, to sprout}

\label{lesson:GU-C218-ugvu}



\begin{quote}
Before the new one: say the Gujarati for whirlwind, then the Gujarati for to spread.
\end{quote}

\subsection*{You'll want to know: \gu{ઊગવું}}
\textbf{\gu{ઊગવું}} — \emph{ūgvũ} — \textquotedblleft{}to rise, to sprout\textquotedblright{}.

\begin{quote}
\textbf{\gu{કાલે} \gu{સૂરજ} \gu{વહેલો} \gu{ઊગશે}.} — \emph{kāle sūraj vahelo ūgśe} — \textquotedblleft{}The sun will rise early tomorrow.\textquotedblright{}
\end{quote}

\begin{cousinweb}
\textbf{\gu{ઊગવું}} (\emph{ūgvũ}) is a tadbhava. It comes from \emph{*udgāti}, \textquotedblleft{}rises, comes forth\textquotedblright{}, a reconstructed form; CDIAL glosses the Gujarati verb \textquotedblleft{}to rise, come out, grow\textquotedblright{} (CDIAL 1954).
\end{cousinweb}

\begin{grammarlens}[title={What you've built}]
One more word for everyday life at A2.
\end{grammarlens}

\subsection*{Your turn}
\begin{itemize}
\raggedright
  \item \emph{Say it:} \emph{ūgvũ}
  \item \emph{Say it:} \emph{ūgvũ}, once more
  \item \emph{Say it:} \emph{kāle sūraj vahelo ūgśe} — The sun will rise early tomorrow.
  \item \emph{Recall:} say the Gujarati for whirlwind, then the Gujarati for to spread, then say \emph{ūgvũ} again
\end{itemize}

\begin{tcolorbox}[breakable,colback=teal!4,colframe=teal!35!black,title={Before you move on}]
What does \gu{ઊગવું} mean? (\textquotedblleft{}To rise, to sprout\textquotedblright{}.) Say it once more.
\end{tcolorbox}
\section[\texorpdfstring{\gu{સૂર્યોદય} (sūryoday)}{સૂર્યોદય (sūryoday)}]{\gu{સૂર્યોદય} (sūryoday) — sunrise}

\label{lesson:GU-C218-suryoday}



\begin{quote}
Before the new one: say the Gujarati for to spread, then the Gujarati for to rise.
\end{quote}

\subsection*{You'll want to know: \gu{સૂર્યોદય}}
\textbf{\gu{સૂર્યોદય}} — \emph{sūryoday} — \textquotedblleft{}sunrise\textquotedblright{}.

\begin{quote}
\textbf{\gu{આવતીકાલે} \gu{સૂર્યોદય} \gu{વહેલો} \gu{થશે}.} — \emph{āvtīkāle sūryoday vahelo thaśe} — \textquotedblleft{}Sunrise will be early tomorrow.\textquotedblright{}
\end{quote}

\begin{cousinweb}
\textbf{\gu{સૂર્યોદય}} (\emph{sūryoday}) is Sanskrit \emph{sūryodaya}, \textquotedblleft{}sunrise\textquotedblright{}, built from \emph{sūrya}, \textquotedblleft{}sun\textquotedblright{}, and \emph{udaya}, \textquotedblleft{}rising\textquotedblright{}, and taken over whole as a learned loan, a tatsama (CDIAL 13574, 1931).
\end{cousinweb}

\begin{grammarlens}[title={What you've built}]
One more word for everyday life at A2.
\end{grammarlens}

\subsection*{Your turn}
\begin{itemize}
\raggedright
  \item \emph{Say it:} \emph{sūryoday}
  \item \emph{Say it:} \emph{sūryoday}, once more
  \item \emph{Say it:} \emph{āvtīkāle sūryoday vahelo thaśe} — Sunrise will be early tomorrow.
  \item \emph{Recall:} say the Gujarati for to spread, then the Gujarati for to rise, then say \emph{sūryoday} again
\end{itemize}

\begin{tcolorbox}[breakable,colback=teal!4,colframe=teal!35!black,title={Before you move on}]
What does \gu{સૂર્યોદય} mean? (\textquotedblleft{}Sunrise\textquotedblright{}.) Say it once more.
\end{tcolorbox}
\section[\texorpdfstring{\gu{આથમવું} (āthamvũ)}{આથમવું (āthamvũ)}]{\gu{આથમવું} (āthamvũ) — to set}

\label{lesson:GU-C218-athamvu}



\begin{quote}
Before the new one: say the Gujarati for to rise, then the Gujarati for sunrise.
\end{quote}

\subsection*{You'll want to know: \gu{આથમવું}}
\textbf{\gu{આથમવું}} — \emph{āthamvũ} — \textquotedblleft{}to set\textquotedblright{}.

\begin{quote}
\textbf{\gu{સાંજે} \gu{સૂરજ} \gu{મોડો} \gu{આથમશે}.} — \emph{sā̃je sūraj moḍo āthamśe} — \textquotedblleft{}The sun will set late in the evening.\textquotedblright{}
\end{quote}

\begin{cousinweb}
\textbf{\gu{આથમવું}} (\emph{āthamvũ}) is a tadbhava. It comes from the Sanskrit phrase \emph{astam eti}, \textquotedblleft{}goes home, sets (of the sun)\textquotedblright{}, built on \emph{asta}, \textquotedblleft{}home\textquotedblright{} (CDIAL 976).
\end{cousinweb}

\begin{grammarlens}[title={What you've built}]
One more word for everyday life at A2.
\end{grammarlens}

\subsection*{Your turn}
\begin{itemize}
\raggedright
  \item \emph{Say it:} \emph{āthamvũ}
  \item \emph{Say it:} \emph{āthamvũ}, once more
  \item \emph{Say it:} \emph{sā̃je sūraj moḍo āthamśe} — The sun will set late in the evening.
  \item \emph{Recall:} say the Gujarati for to rise, then the Gujarati for sunrise, then say \emph{āthamvũ} again
\end{itemize}

\begin{tcolorbox}[breakable,colback=teal!4,colframe=teal!35!black,title={Before you move on}]
What does \gu{આથમવું} mean? (\textquotedblleft{}To set\textquotedblright{}.) Say it once more.
\end{tcolorbox}
\section[\texorpdfstring{\gu{સૂર્યાસ્ત} (sūryāst)}{સૂર્યાસ્ત (sūryāst)}]{\gu{સૂર્યાસ્ત} (sūryāst) — sunset}

\label{lesson:GU-C218-suryast}



\begin{quote}
Before the new one: say the Gujarati for sunrise, then the Gujarati for to set.
\end{quote}

\subsection*{You'll want to know: \gu{સૂર્યાસ્ત}}
\textbf{\gu{સૂર્યાસ્ત}} — \emph{sūryāst} — \textquotedblleft{}sunset\textquotedblright{}.

\begin{quote}
\textbf{\gu{આજે} \gu{સૂર્યાસ્ત} \gu{સુંદર} \gu{હશે}.} — \emph{āje sūryāst sundar haśe} — \textquotedblleft{}The sunset will be beautiful today.\textquotedblright{}
\end{quote}

A masculine noun.

\begin{cousinweb}
\textbf{\gu{સૂર્યાસ્ત}} (\emph{sūryāst}) is Sanskrit \emph{sūryāsta}, \textquotedblleft{}sunset\textquotedblright{}, built from \emph{sūrya}, \textquotedblleft{}sun\textquotedblright{}, and \emph{asta}, \textquotedblleft{}setting\textquotedblright{} (CDIAL 973), taken over whole as a tatsama.
\end{cousinweb}

\begin{grammarlens}[title={What you've built}]
One more word for everyday life at A2.
\end{grammarlens}

\subsection*{Your turn}
\begin{itemize}
\raggedright
  \item \emph{Say it:} \emph{sūryāst}
  \item \emph{Say it:} \emph{sūryāst}, once more
  \item \emph{Say it:} \emph{āje sūryāst sundar haśe} — The sunset will be beautiful today.
  \item \emph{Recall:} say the Gujarati for sunrise, then the Gujarati for to set, then say \emph{sūryāst} again
\end{itemize}

\begin{tcolorbox}[breakable,colback=teal!4,colframe=teal!35!black,title={Before you move on}]
What does \gu{સૂર્યાસ્ત} mean? (\textquotedblleft{}Sunset\textquotedblright{}.) Say it once more.
\end{tcolorbox}
\section[\texorpdfstring{\gu{અજવાળું} (ajvāḷũ)}{અજવાળું (ajvāḷũ)}]{\gu{અજવાળું} (ajvāḷũ) — daylight, light}

\label{lesson:GU-C218-ajvalu}



\begin{quote}
Before the new one: say the Gujarati for to set, then the Gujarati for sunset.
\end{quote}

\subsection*{You'll want to know: \gu{અજવાળું}}
\textbf{\gu{અજવાળું}} — \emph{ajvāḷũ} — \textquotedblleft{}daylight, light\textquotedblright{}.

\begin{quote}
\textbf{\gu{પરોઢે} \gu{અજવાળું} \gu{થશે}.} — \emph{paroḍhe ajvāḷũ thaśe} — \textquotedblleft{}It will get light at dawn.\textquotedblright{}
\end{quote}

A neuter noun.

\begin{cousinweb}
\textbf{\gu{અજવાળું}} (\emph{ajvāḷũ}) is usually linked to Sanskrit \emph{ujjvala}, \textquotedblleft{}blazing up, bright\textquotedblright{}. CDIAL lists the related Gujarati \emph{ujḷũ}, \textquotedblleft{}bright\textquotedblright{}, there, but not \emph{ajvāḷũ} itself, so the exact path is not settled (CDIAL 1670).
\end{cousinweb}

\begin{grammarlens}[title={What you've built}]
One more word for everyday life at A2.
\end{grammarlens}

\subsection*{Your turn}
\begin{itemize}
\raggedright
  \item \emph{Say it:} \emph{ajvāḷũ}
  \item \emph{Say it:} \emph{ajvāḷũ}, once more
  \item \emph{Say it:} \emph{paroḍhe ajvāḷũ thaśe} — It will get light at dawn.
  \item \emph{Recall:} say the Gujarati for to set, then the Gujarati for sunset, then say \emph{ajvāḷũ} again
\end{itemize}

\begin{tcolorbox}[breakable,colback=teal!4,colframe=teal!35!black,title={Before you move on}]
What does \gu{અજવાળું} mean? (\textquotedblleft{}Daylight, light\textquotedblright{}.) Say it once more.
\end{tcolorbox}
\section[(dialogue)]{Weather and the seasons - a review}

\label{lesson:GU-R218-review-213-218}



\begin{quote}
Say the Gujarati for weather and sunshine.
\end{quote}

\subsection*{Your turn}
Say these aloud:

\begin{itemize}
\raggedright
  \item weather, sunshine, cold, cloudy, to pour down
  \item to fall, pouring, a drop, to drip, to get soaked
  \item to freeze, frost, dew, thick, lukewarm
  \item to blaze, sultry heat, humidity, temperature, degree
  \item air, to blow, dust storm, whirlwind, to spread
  \item to rise, sunrise, to set, sunset, daylight
\end{itemize}

\begin{tcolorbox}[breakable,colback=teal!4,colframe=teal!35!black,title={Before you move on}]
Weather? (\textbf{\gu{હવામાન}}.) Sunshine? (\textbf{\gu{તડકો}}.) Cold? (\textbf{\gu{ઠંડું}}.) Cloudy? (\textbf{\gu{વાદળછાયું}}.) To pour down? (\textbf{\gu{વરસવું}}.) To fall? (\textbf{\gu{પડવું}}.) Pouring? (\textbf{\gu{ધોધમાર}}.) A drop? (\textbf{\gu{ટીપું}}.) To drip? (\textbf{\gu{ટપકવું}}.) To get soaked? (\textbf{\gu{પલળવું}}.) To freeze? (\textbf{\gu{ઠરવું}}.) Frost? (\textbf{\gu{હિમ}}.) Dew? (\textbf{\gu{ઓસ}}.) Thick? (\textbf{\gu{ગાઢ}}.) Lukewarm? (\textbf{\gu{હૂંફાળું}}.) To blaze? (\textbf{\gu{તપવું}}.) Sultry heat? (\textbf{\gu{ઉકળાટ}}.) Humidity? (\textbf{\gu{ભેજ}}.) Temperature? (\textbf{\gu{તાપમાન}}.) Degree? (\textbf{\gu{ડિગ્રી}}.) Air? (\textbf{\gu{હવા}}.) To blow? (\textbf{\gu{ફૂંકાવું}}.) Dust storm? (\textbf{\gu{આંધી}}.) Whirlwind? (\textbf{\gu{વંટોળ}}.) To spread? (\textbf{\gu{ફેલાવું}}.) To rise? (\textbf{\gu{ઊગવું}}.) Sunrise? (\textbf{\gu{સૂર્યોદય}}.) To set? (\textbf{\gu{આથમવું}}.) Sunset? (\textbf{\gu{સૂર્યાસ્ત}}.) Daylight? (\textbf{\gu{અજવાળું}}.)
\end{tcolorbox}
